\documentclass[10pt,a4paper]{amsart}
\usepackage[margin=19mm]{geometry}
\usepackage{amsmath,amssymb,mathtools}
\usepackage[scr=rsfs,cal=euler]{mathalfa}
\usepackage{xfrac}
\usepackage[unicode,hidelinks,bookmarks=false]{hyperref}
\newcommand{\ie}{i.e.\ }
\newcommand{\cf}{cf.\ }
\newcommand{\resp}{respectively\ }
\newcommand{\Subsection}{\subsection}
\providecommand{\nobreakdash}{\nobreak\hbox{-}}
\setlength{\emergencystretch}{2em}
\multlinegap0pt
\usepackage{bm}
\usepackage{bbm}
\usepackage{dsfont}
\usepackage{xspace}
\usepackage{cancel}
\usepackage{mathtools}
\usepackage{amsthm}
\makeatletter
\@ifundefined{thm}{\newtheorem{thm}{Thm}}{}
\@ifundefined{lem}{\newtheorem{lem}[thm]{Lemma}}{}
\makeatother
\title[On the density version of quadratic Waring's problem and the]{On the density version of quadratic Waring's problem and the quadratic Waring--Goldbach problem}
\author{Zi Li Lim}
\date{Source: arXiv:2508.14939v3 (CC BY 4.0)}
\begin{document}
\maketitle

\noindent\emph{Source and licence.} On the density version of quadratic Waring's problem and the quadratic Waring--Goldbach problem. Version arXiv:2508.14939v3 (CC BY 4.0), distributed under the Creative Commons Attribution 4.0 International licence, as recorded in the arXiv licence metadata for this exact version and shown on the abstract page https://arxiv.org/abs/2508.14939v3. Full rights notice: licenses/arxiv-2508.14939-CC-BY-4.0.txt.\par\smallskip

\noindent\textit{Editorial scope.} Counted unit: the Lem whose source label is \texttt{lemma for base cases} in the article (source0.tex lines 972--981, source label \texttt{lemma for base cases}), delivered together with the article's own complete original proof of it (source0.tex lines 983--1079, one \texttt{proof} environment of 97 lines). No mathematics is written, repaired or re-derived: statement and proof are the article's own text line for line. Required context retained uncounted: combinatorial lemma (lines 847--856). Every definition, display and label of those lines is retained unchanged. Omitted from this package and not redistributed: 1--846, 857--971, 982--982, 1080--1434. Nothing inside a retained excerpt is omitted or altered: every retained body is delivered exactly as the article's lines are written, so no omission receipt of any kind accompanies this package. Citations: the retained excerpts carry no citation command at all, so no bibliography entry of the article is transcribed and no reference is redistributed. Reference adaptation: none was needed and none was made; every reference inside the delivered unit resolves inside it, so no unresolved reference or citation is produced. Printed slips kept literal: no printed word, symbol, hypothesis or number of the counted statement or of its proof was altered. Layout and class: the article's own class is \texttt{amsart} with packages including amsmath, amssymb, amsthm, amsfonts, mathrsfs, mathtools, xcolor, xy, geometry, biblatex, hyperref, xfrac; the wrapper is the lane's pinned \texttt{amsart} editorial wrapper at 10pt on A4, which restates the theorem-like environments the retained text needs and adds the article's own macro definitions that the retained text needs (none), with extra wrapper packages (bm, bbm, dsfont, xspace, cancel, mathtools). The delivered numbering is the unit's own. Assets: no retained span contains a figure environment, a graphics-inclusion command, a PDF-page inclusion, a table or a TikZ picture; no image file is copied into this package, the package declares no asset dependency and no image file is redistributed with it. Licence: version arXiv:2508.14939v3 of this article is distributed under the Creative Commons Attribution 4.0 International licence, as recorded in the arXiv licence metadata for this exact version and shown on the abstract page https://arxiv.org/abs/2508.14939v3. A full rights notice is delivered at licenses/arxiv-2508.14939-CC-BY-4.0.txt. Characters: every retained excerpt is pure ASCII, so the delivered engine, XeLaTeX with its default Latin Modern text font, reports no missing character.
\begin{lem}\label{combinatorial lemma}
    Let $s\geq 6, n\geq 3$ or $s=5, n\geq 5$. For each integer $1\leq j\leq s$, let $\{a_{i,j}\}_{0\leq i<n}$ be a decreasing sequence in interval $[0,1]$. Suppose 
    \begin{equation*}
        \text{$\frac{1}{ns}\sum_{\substack{0\leq i<n\\1\leq j\leq s}}a_{i,j}>d_s$ and $a_{0,j}\neq 0$ for all $1\leq j\leq s$.}
    \end{equation*}
    Then, there exist indices $0\leq i_1,i_2,...,i_s<n$ with $i_1+i_2+...+i_s\geq 2n$ such that
    \begin{equation*}
        \text{$\frac{1}{s}\sum_{1\leq j\leq s}a_{i_j,j}>d_s$ and $a_{i_j,j}\neq 0$ for all $1\leq j\leq s$.}
    \end{equation*}
\end{lem}

\begin{lem}\label{lemma for base cases}
    Let $s\geq 12$ be an integer, and $D\geq \frac{1}{2}$ be a real number. For each integer $1\leq j\leq s$, let $\{a_{i,j}\}_{0\leq i<3}$ be a decreasing sequence in interval $[0,1]$. Suppose 
    \begin{equation*}
        \text{$\frac{1}{3s}\sum_{\substack{0\leq i<3\\1\leq j\leq s}}a_{i,j}>D$ and $a_{0,j}\neq 0$ for all $1\leq j\leq s$.}
    \end{equation*}
    Then, there exist indices $0\leq i_1,i_2,...,i_s<3$ with $i_1+i_2+...+i_s\geq 6$ such that
    \begin{equation*}
        \text{$\frac{1}{s}\sum_{1\leq j\leq s}a_{i_j,j}>D$ and $a_{i_j,j}\neq 0$ for all $1\leq j\leq s$.}
    \end{equation*}
\end{lem}

\begin{proof}[Proof of Lemma \ref{lemma for base cases} and bases cases of Lemma \ref{combinatorial lemma} when $s\geq 12$]
    Let $R$ be the number of nonzero elements among $a_{2,j}$ where $1\leq j\leq s$. We will prove Lemma \ref{lemma for base cases} via cases-by-cases analysis based on the values of $R$. Note that the number of nonzero elements among $a_{i,j}$ where $i=1,2$ and $1\leq j\leq s$ is $\geq 7$, because this number is greater than $s/2$.\\

    \noindent\underline{\textbf{Case 0:}} $R=0$

    There are at least $7$ nonzero elements among $a_{i,j}$ where $i=1,2$ and $1\leq j\leq s$. By relabeling, we may assume $a_{1,j}\neq 0$ for $1\leq j\leq 7$. Suppose Lemma \ref{lemma for base cases} is false, then
    \begin{align*}
        &(a_{1,1}+a_{1,2}+a_{1,3}+a_{1,4}+a_{1,5}+a_{1,6})+(a_{0,7}+a_{0,8}+...+a_{0,s})\leq sD,\\
        &a_{0,1}+(a_{1,2}+a_{1,3}+a_{1,4}+a_{1,5}+a_{1,6}+a_{1,7})+(a_{0,8}+a_{0,9}+...+a_{0,s})\leq sD.
    \end{align*}
    Summing the inequalities above and using the facts that $a_{i,j}$ is decreasing in $i$ and $R=0$, we have
    \begin{equation*}
        \sum_{\substack{0\leq i<3\\1\leq j\leq s}}a_{i,j}-(a_{0,2}+a_{0,3}+a_{0,4}+a_{0,5}+a_{0,6})\leq 2sD.
    \end{equation*}
    This implies $3sD-5<2sD$, which leads to a contradiction.\\

    The strategies for the remaining cases are similar to Case $0$.\\

    \noindent\underline{\textbf{Case 1:}} $R=1$

    There are at least $7$ nonzero elements among $a_{i,j}$ where $i=1,2$ and $1\leq j\leq s$. By relabeling, we may assume $a_{1,j}\neq 0$ for $1\leq j\leq 6$ and $a_{2,1}\neq 0$. Suppose Lemma \ref{lemma for base cases} is false, then
    \begin{align*}
        &a_{2,1}+(a_{1,2}+a_{1,3}+a_{1,4}+a_{1,5})+(a_{0,6}+a_{0,7}+...+a_{0,s})\leq sD,\\
        &(a_{1,1}+a_{1,2}+a_{1,3}+a_{1,4}+a_{1,5}+a_{1,6})+(a_{0,7}+a_{0,8}+...+a_{0,s})\leq sD.
    \end{align*}
    Summing the inequalities above and using the facts that $a_{i,j}$ is decreasing in $i$ and $R=1$, we have
    \begin{equation*}
        \sum_{\substack{0\leq i<3\\1\leq j\leq s}}a_{i,j}-(a_{0,1}+a_{0,2}+a_{0,3}+a_{0,4}+a_{0,5})\leq 2sD.
    \end{equation*}
    This implies $3sD-5<2sD$, which leads to a contradiction.\\

    \noindent\underline{\textbf{Case 2:}} $R=2$

    There are at least $7$ nonzero elements among $a_{i,j}$ where $i=1,2$ and $1\leq j\leq s$. By relabeling, we may assume $a_{1,j}\neq 0$ for $1\leq j\leq 5$ and $a_{2,1},a_{2,2}\neq 0$. Suppose Lemma \ref{lemma for base cases} is false, then
    \begin{align*}
        &(a_{2,1}+a_{2,2})+(a_{1,3}+a_{1,4})+(a_{0,5}+a_{0,6}+...+a_{0,s})\leq sD,\\
        &a_{2,1}+(a_{1,2}+a_{1,3}+a_{1,4}+a_{1,5})+(a_{0,6}+a_{0,7}+...+a_{0,s})\leq sD.
    \end{align*}
    Summing the inequalities above and using the facts that $a_{i,j}$ is decreasing in $i$ and $R=2$, we have
    \begin{equation*}
        \sum_{\substack{0\leq i<3\\1\leq j\leq s}}a_{i,j}-(a_{0,1}+a_{1,1}+a_{0,2}+a_{0,3}+a_{0,4})\leq 2sD.
    \end{equation*}
    This implies $3sD-5<2sD$, which leads to a contradiction.\\

    \noindent\underline{\textbf{Case 3:}} $R=3$

    There are at least $7$ nonzero elements among $a_{i,j}$ where $i=1,2$ and $1\leq j\leq s$. By relabeling, we may assume $a_{1,j}\neq 0$ for $1\leq j\leq 4$ and $a_{2,1},a_{2,2},a_{2,3}\neq 0$. Suppose Lemma \ref{lemma for base cases} is false, then
    \begin{align*}
        &(a_{2,1}+a_{2,2}+a_{2,3})+(a_{0,4}+a_{0,5}+a_{0,6}+a_{0,7}+...+a_{0,s})\leq sD,\\
        &(a_{2,1}+a_{2,2})+(a_{1,3}+a_{1,4})+(a_{0,5}+a_{0,6}+...+a_{0,s})\leq sD.
    \end{align*}
    Summing the inequalities above and using the facts that $a_{i,j}$ is decreasing in $i$ and $R=3$, we have
    \begin{equation*}
        \sum_{\substack{0\leq i<3\\1\leq j\leq s}}a_{i,j}-(a_{0,1}+a_{1,1}+a_{0,2}+a_{1,2}+a_{0,3})\leq 2sD.
    \end{equation*}
    This implies $3sD-5<2sD$, which leads to a contradiction.\\

    \noindent\underline{\textbf{Case 4:}} $R=4$

    By relabeling, we may assume $a_{i,j}\neq 0$ for $1\leq i\leq 2,1\leq j\leq 4$. Suppose Lemma \ref{lemma for base cases} is false, then
    \begin{align*}
        &(a_{2,1}+a_{2,2})+(a_{1,3}+a_{1,4})+(a_{0,5}+a_{0,6}+...+a_{0,s})\leq sD,\\
        &(a_{1,1}+a_{1,2})+(a_{2,3}+a_{2,4})+(a_{0,5}+a_{0,6}+...+a_{0,s})\leq sD.
    \end{align*}
    Summing the inequalities above and using the facts that $a_{i,j}$ is decreasing in $i$ and $R=4$, we have
    \begin{equation*}
        \sum_{\substack{0\leq i<3\\1\leq j\leq s}}a_{i,j}-(a_{0,1}+a_{0,2}+a_{0,3}+a_{0,4})\leq 2sD.
    \end{equation*}
    This implies $3sD-4<2sD$, which leads to a contradiction.\\

    \noindent\underline{\textbf{Case 5:}} $R=5$

    By relabeling, we may assume $a_{i,j}\neq 0$ for $1\leq i\leq 2,1\leq j\leq 5$. Suppose Lemma \ref{lemma for base cases} is false, then
    \begin{align*}
        &(a_{2,1}+a_{2,2})+(a_{1,3}+a_{1,4}+a_{1,5})+(a_{0,6}+a_{0,7}+...+a_{0,s})\leq sD,\\
        &(a_{1,1}+a_{1,2})+(a_{2,3}+a_{2,4}+a_{2,5})+(a_{0,6}+a_{0,7}+...+a_{0,s})\leq sD.
    \end{align*}
    Summing the inequalities above and using the facts that $a_{i,j}$ is decreasing in $i$ and $R=5$, we have
    \begin{equation*}
        \sum_{\substack{0\leq i<3\\1\leq j\leq s}}a_{i,j}-(a_{0,1}+a_{0,2}+a_{0,3}+a_{0,4}+a_{0,5})\leq 2sD.
    \end{equation*}
    This implies $3sD-5<2sD$, which leads to a contradiction.\\

    \noindent\underline{\textbf{Case 6:}} $R\geq 6$

    By relabeling, we may assume $a_{i,j}\neq 0$ for $1\leq i\leq 2,1\leq j\leq 6$. Suppose Lemma \ref{lemma for base cases} is false, then
    \begin{align*}
        &(a_{0,1}+a_{0,2}+a_{0,3})+(a_{2,4}+a_{2,5}+a_{2,6})+(a_{0,7}+a_{0,8}+...+a_{0,s})\leq sD,\\
        &(a_{1,1}+a_{1,2}+a_{1,3}+a_{1,4}+a_{1,5}+a_{1,6})+(a_{0,7}+a_{0,8}+...+a_{0,s})\leq sD,\\
        &(a_{2,1}+a_{2,2}+a_{2,3})+(a_{0,4}+a_{0,5}+a_{0,6})+(a_{0,7}+a_{0,8}+...+a_{0,s})\leq sD.
    \end{align*}
    Summing the inequalities above and using the facts that $a_{i,j}$ is decreasing in $i$, we have
    \begin{equation*}
        \sum_{\substack{0\leq i<3\\1\leq j\leq s}}a_{i,j}\leq 3sD.
    \end{equation*}
    This implies $3sD<3sD$, which leads to a contradiction.\\
\end{proof}

\end{document}
